%%=============================================================================
%% 第 3 章 方法
%%=============================================================================

\chapter{基于表格型 Q-Learning 的时间窗优化方法}
\label{chap:method}

% 引言
在上一章中，本研究通过文献综述发现：现有强化学习方法（如 MARS）虽在 BCI 特征选择中表现优异，但其复杂的多智能体设置和深度网络结构带来了巨大的计算开销。基于此，本章提出一种基于表格型 Q-Learning 的运动想象时间窗优化方法，旨在以简洁的表格型强化学习框架实现自适应时间窗选择，在保持性能的同时显著降低计算复杂度。

本章首先将时间窗优化问题形式化为马尔可夫决策过程（第\ref{sec:formulation}节），明确定义状态空间、动作空间、奖励函数和状态转移模型；其次，详细阐述表格型 Q-Learning 算法的核心原理与具体设计（第\ref{sec:qlearning}节），包括$\varepsilon$-贪婪探索策略、Q 值更新规则和推理过程；再次，从理论上分析算法的收敛性保证（第\ref{sec:convergence}节），证明本研究方法在有限状态空间下能够收敛到最优时间窗选择策略；最后，系统对比表格型 Q-Learning 与深度 Q 网络（DQN）的理论特性（第\ref{sec:comparison}节），阐明本研究选择表格型方法的理论依据与适用边界。

\section{问题形式化}
\label{sec:formulation}

本节将时间窗优化问题形式化为马尔可夫决策过程（Markov Decision Process, MDP），并定义优化目标。

\subsection{马尔可夫决策过程建模}
\label{subsec:mdp}

本研究将时间窗优化问题建模为马尔可夫决策过程，形式化定义为五元组 $(\mathcal{S}, \mathcal{A}, \mathcal{R}, \mathcal{P}, \gamma)$。各组成部分定义如下：

\textbf{状态空间 $\mathcal{S}$}。状态定义为离散化的时间窗参数对 $(t_{\text{start}}, t_{\text{end}})$，其中 $t_{\text{start}}$ 表示时间窗起始时刻，$t_{\text{end}}$ 表示时间窗终止时刻。时间参数离散化步长为 $\Delta t = 0.2$ 秒，为保证时间窗具有最小有效长度，约束条件为 $t_{\text{end}} - t_{\text{start}} \geq 0.5$ 秒。状态空间可形式化表示为：
\begin{equation}
\mathcal{S} = \left\{ (t_s, t_e) \mid t_s \in \mathcal{T}_s, t_e \in \mathcal{T}_e, t_e - t_s \geq 0.5 \right\}
\end{equation}
其中，$\mathcal{T}_s = \{0.2k \mid k \in \mathbb{Z}, 0 \leq k \leq 15\}$，$\mathcal{T}_e = \{0.2k \mid k \in \mathbb{Z}, 5 \leq k \leq 20\}$。

状态空间的大小可通过以下公式计算：
\begin{equation}
|\mathcal{S}| = \sum_{k=0}^{15} \max\left(0, \left\lfloor \frac{4.0 - \max(1.0, 0.2k + 0.5)}{0.2} \right\rfloor - 5 + 1\right) = 136
\end{equation}
即有效状态数量为 136 个。

\textbf{动作空间 $\mathcal{A}$}。动作空间定义为离散的时间窗边界调整操作：
\begin{equation}
\mathcal{A} = \{a_0, a_1, a_2, a_3\}
\end{equation}
各动作的语义定义为：$a_0$ 表示将 $t_{\text{start}}$ 提前 $\Delta t$（左边界左移，时间窗变宽），$a_1$ 表示将 $t_{\text{start}}$ 延后 $\Delta t$（左边界右移，时间窗变窄），$a_2$ 表示将 $t_{\text{end}}$ 提前 $\Delta t$（右边界左移，时间窗变窄），$a_3$ 表示将 $t_{\text{end}}$ 延后 $\Delta t$（右边界右移，时间窗变宽）。动作执行后若导致时间窗超出有效范围（$t_{\text{start}} < 0$、$t_{\text{end}} > 4.0$ 或 $t_{\text{end}} - t_{\text{start}} < 0.5$），则状态保持不变。值得注意的是，本任务中的状态转移是确定性的，这是时间窗参数离散化带来的便利特性；在连续参数场景中，转移可能具有随机性。

\textbf{奖励函数 $\mathcal{R}$}。奖励函数定义为状态 - 动作对 $(s, a)$ 的即时反馈信号，形式化表示为：
\begin{equation}
\mathcal{R}: \mathcal{S} \times \mathcal{A} \rightarrow \mathbb{R}, \quad r(s, a) = \text{Accuracy}_{\text{CV}}(s'; \mathcal{D}_{\text{train}})
\end{equation}
其中，$s' = f(s, a)$ 为动作执行后的新状态，$\text{Accuracy}_{\text{CV}}$ 表示在新状态对应的时间窗下，基础分类器在训练集 $\mathcal{D}_{\text{train}}$ 上的 $k$ 折交叉验证准确率。本研究采用 CSP-SVM 作为基础分类器\cite{blankertz2008optimizing}，使用 5 折交叉验证计算奖励信号。

\textbf{状态转移 $\mathcal{P}$}。状态转移为确定性过程，由状态转移函数 $f: \mathcal{S} \times \mathcal{A} \rightarrow \mathcal{S}$ 唯一确定：
\begin{equation}
\mathcal{P}(s' \mid s, a) =
\begin{cases}
1, & \text{if } s' = f(s, a) \\
0, & \text{otherwise}
\end{cases}
\end{equation}
确定性转移假设在本任务中成立，因为时间窗参数的调整是可预测的离散操作。

\textbf{折扣因子 $\gamma$}。折扣因子 $\gamma \in [0, 1)$ 用于平衡即时奖励与长期奖励的权重。本研究设 $\gamma = 0.95$，使智能体在关注即时分类性能的同时，兼顾长期累积奖励。

\subsection{优化目标}
\label{subsec:objective}

时间窗优化的目标是找到最优策略 $\pi^*$，使得期望累积奖励最大化：
\begin{equation}
\pi^* = \arg\max_{\pi} \mathbb{E}_{\pi} \left[ \sum_{t=0}^{\infty} \gamma^t r_t \right]
\end{equation}
其中，$r_t$ 为时刻 $t$ 的即时奖励。在本任务中，由于状态转移是确定性的且 episode 长度有限，优化目标等价于找到最优时间窗 $(t_{\text{start}}^*, t_{\text{end}}^*)$，使得分类准确率最高。

\section{表格型 Q-Learning 算法}
\label{sec:qlearning}

\subsection{算法原理}
\label{subsec:algorithm}

Q-Learning 是一种经典的无模型强化学习算法，由 Watkins 和 Dayan 于 1992 年提出\cite{watkins1992q}。其核心思想是通过维护状态 - 动作价值表（Q 表），迭代估计最优动作价值函数 $Q^*(s, a)$。

\textbf{Q 函数的定义}。动作价值函数 $Q(s, a)$ 表示在状态 $s$ 执行动作 $a$ 后，遵循当前策略可获得的期望累积奖励：
\begin{equation}
Q(s, a) = \mathbb{E} \left[ \sum_{t=0}^{\infty} \gamma^t r_t \mid s_0 = s, a_0 = a \right]
\end{equation}
其中，$\gamma \in [0, 1)$ 为折扣因子，用于平衡即时奖励与长期奖励的权重。最优动作价值函数 $Q^*(s, a)$ 满足 Bellman 最优方程：
\begin{equation}
Q^*(s, a) = \mathbb{E}_{s' \sim \mathcal{P}} \left[ r + \gamma \max_{a'} Q^*(s', a') \mid s, a \right]
\end{equation}

\textbf{时序差分学习}。Q-Learning 采用时序差分（Temporal Difference, TD）学习方法，通过自举（bootstrapping）从后续状态的估计值中更新当前 Q 值。Q 值更新遵循以下规则：
\begin{equation}
Q(s_t, a_t) \leftarrow Q(s_t, a_t) + \alpha \left[ r_{t+1} + \gamma \max_{a} Q(s_{t+1}, a) - Q(s_t, a_t) \right]
\end{equation}
其中，$\alpha \in (0, 1]$ 为学习率，控制新信息覆盖旧信息的程度；$r_{t+1} + \gamma \max_{a} Q(s_{t+1}, a)$ 为 TD 目标，表示基于经验估计的真实 Q 值；方括号内的差值称为\textbf{TD 误差}，表示当前估计值与目标值之间的差距。

% \textbf{探索与利用的权衡}。Q-Learning 采用$\varepsilon$-贪婪策略平衡探索与利用：以概率 $1-\varepsilon$ 选择当前 Q 值最大的动作（利用），以概率 $\varepsilon$ 随机选择动作（探索）\cite{sutton2018reinforcement}。训练初期设置较高的$\varepsilon$ 值鼓励充分探索状态空间；随训练进行，$\varepsilon$ 按指数衰减，逐步转向利用已知的高 Q 值动作。

% \textbf{收敛性保证}。Watkins 和 Dayan（1992）证明了 Q-Learning 在有限马尔可夫决策过程下的收敛性\cite{watkins1992q}：若所有状态 - 动作对被无限次访问，且学习率序列满足 Robbins-Monro 条件，则 Q 值序列以概率 1 收敛到最优动作价值函数 $Q^*$。

\subsection{算法设计}
\label{subsec:design}

本研究提出的表格型 Q-Learning 时间窗优化算法如算法~\ref{alg:qlearning} 所示。

\begin{algorithm}[H]
\caption{表格型 Q-Learning 时间窗优化算法}
\label{alg:qlearning}
\begin{algorithmic}[1]
\REQUIRE EEG 训练数据 $\mathcal{D} = \{(\mathbf{X}_i, y_i)\}_{i=1}^N$
\REQUIRE 状态空间 $\mathcal{S}$（136 个离散时间窗状态）
\REQUIRE 动作空间 $\mathcal{A} = \{0, 1, 2, 3\}$
\REQUIRE 学习率 $\alpha = 0.2$，折扣因子 $\gamma = 0.95$
\REQUIRE 初始探索率 $\varepsilon_{\text{init}} = 0.9$，最小探索率 $\varepsilon_{\text{min}} = 0.05$，探索率衰减 $\varepsilon_{\text{decay}} = 0.995$
\REQUIRE 训练轮数 $n_{\text{episodes}} = 100$，每轮最大步数 $\text{max\_steps} = 40$
\ENSURE 最优时间窗 $(t_{\text{start}}^*, t_{\text{end}}^*)$，Q 值表 $Q$

\STATE 初始化 Q 表：$Q[s, a] \leftarrow 0, \forall s \in \mathcal{S}, a \in \mathcal{A}$
\STATE $\varepsilon \leftarrow \varepsilon_{\text{init}}$
\FOR{$\text{episode} = 1$ \TO $n_{\text{episodes}}$}
    \STATE 随机初始化状态 $s_0 = (t_{\text{start}}, t_{\text{end}})$
    \FOR{$\text{step} = 1$ \TO $\text{max\_steps}$}
        \STATE \textbf{if} $\text{random}() < \varepsilon$ \textbf{then}
            \STATE $a_t \leftarrow$ 随机选择动作 $\in \mathcal{A}$
        \STATE \textbf{else}
            \STATE $a_t \leftarrow \arg\max_a Q[s_t, a]$
        \STATE \textbf{end if}
        \STATE 执行动作，状态转移：$s_{t+1} \leftarrow f(s_t, a_t)$
        \STATE 计算奖励：$r_t \leftarrow \text{EvaluateWindow}(s_{t+1}, \mathcal{D})$
        \STATE Q 值更新：$Q[s_t, a_t] \leftarrow Q[s_t, a_t] + \alpha[r_t + \gamma\max_a Q[s_{t+1}, a] - Q[s_t, a_t]]$
        \STATE $s_t \leftarrow s_{t+1}$
    \ENDFOR
    \STATE 衰减探索率：$\varepsilon \leftarrow \max(\varepsilon_{\text{min}}, \varepsilon \times \varepsilon_{\text{decay}})$
\ENDFOR
\STATE $s^* \leftarrow \arg\max_s \max_a Q[s, a]$
\STATE \RETURN $(t_{\text{start}}^*, t_{\text{end}}^*) = s^*$, $Q$
\end{algorithmic}
\end{algorithm}

奖励计算过程如算法~\ref{alg:evaluate} 所示。

\begin{algorithm}[H]
\caption{时间窗评估函数 $\text{EvaluateWindow}$}
\label{alg:evaluate}
\begin{algorithmic}[1]
\REQUIRE 状态 $s = (t_{\text{start}}, t_{\text{end}})$，EEG 训练数据 $\mathcal{D} = \{(\mathbf{X}_i, y_i)\}_{i=1}^N$
\REQUIRE CSP 分量数 $m = 4$，SVM 正则化参数 $C = 1.0$，交叉验证折数 $k = 5$
\ENSURE 交叉验证准确率 $\text{acc}$

\STATE \textit{// 截取时间窗}
\FOR{$i = 1$ \TO $N$}
    \STATE $\mathbf{X}_i^{\text{window}} \leftarrow \mathbf{X}_i[:, t_{\text{start}}:t_{\text{end}}]$
\ENDFOR

\STATE \textit{// CSP 特征提取}
\STATE 计算各类别协方差矩阵：$\mathbf{C}_c = \frac{1}{N_c} \sum_{i \in \text{class}_c} \mathbf{X}_i^{\text{window}} (\mathbf{X}_i^{\text{window}})^\top$
\STATE 求解广义特征值问题：$\mathbf{C}_1 \mathbf{w} = \lambda (\mathbf{C}_1 + \mathbf{C}_2) \mathbf{w}$
\STATE 构建空间滤波器：$\mathbf{W} = [\mathbf{w}_1, \dots, \mathbf{w}_m]^\top$
\STATE 提取对数方差特征：$\mathbf{f}_i = \log(\text{var}(\mathbf{W} \mathbf{X}_i^{\text{window}}))$

\STATE \textit{// Z-score 标准化}
\STATE $\mathbf{f}_i \leftarrow (\mathbf{f}_i - \mu_f) / \sigma_f$

\STATE \textit{// $k$ 折交叉验证}
\FOR{$\text{fold} = 1$ \TO $k$}
    \STATE 划分训练集 $\mathcal{D}_{\text{train}}^{\text{fold}}$ 和验证集 $\mathcal{D}_{\text{val}}^{\text{fold}}$
    \STATE 训练线性 SVM：$\min_{\mathbf{w},b} \frac{1}{2}\|\mathbf{w}\|^2 + C \sum_i \xi_i$
    \STATE 计算验证集准确率 $\text{acc}_{\text{fold}}$
\ENDFOR

\STATE $\text{acc} \leftarrow \frac{1}{k} \sum_{\text{fold}=1}^k \text{acc}_{\text{fold}}$
\STATE \RETURN $\text{acc}$
\end{algorithmic}
\end{algorithm}

\textbf{动作选择策略}。采用$\varepsilon$-贪婪策略平衡探索与利用：训练初期以高概率（$\varepsilon = 0.9$）随机选择动作，鼓励充分探索状态空间；随训练进行，$\varepsilon$ 按指数衰减（衰减率 0.995），逐步转向利用已知的高 Q 值动作；最终保持最小探索率（$\varepsilon = 0.05$），避免陷入局部最优。参数选择依据如下：学习率$\alpha = 0.2$参考了经典 Q-Learning 文献的常用设置；折扣因子$\gamma = 0.95$使智能体在关注即时性能的同时兼顾长期奖励；训练轮数$n_{\text{episodes}} = 100$的设定基于状态空间规模（136 个状态），100 轮 × 40 步 = 4000 步的总探索次数足以使 544 个状态 - 动作对被充分访问（平均每对约 7.4 次）。

\textbf{奖励计算}。核心步骤包括时间窗截取、CSP 特征提取和 SVM 交叉验证。需要说明的是，每次状态转移都需要进行一次完整的 5 折交叉验证，这在训练阶段（100 episodes × 40 steps = 4000 次评估）会产生较大的计算开销。然而，这属于\textbf{离线训练阶段}的一次性开销，而\textbf{在线推理阶段}仅需查表操作（时间复杂度$O(1)$），完全满足在线 BCI 系统的实时性要求。在工程实现中，可采用缓存机制避免重复计算相同时间窗的奖励值，进一步提升训练效率。


\subsection{推理过程}
\label{subsec:inference}

训练完成后，最优时间窗通过查表获得：
\begin{equation}
s^* = \arg\max_{s \in \mathcal{S}} \max_{a \in \mathcal{A}} Q(s, a)
\end{equation}
推理过程仅需遍历 Q 表，找到最大 Q 值对应的状态，时间复杂度为 $O(|\mathcal{S}| \cdot |\mathcal{A}|) = O(136 \times 4) = O(544)$，在实际计算中可忽略不计。

\section{收敛性分析}
\label{sec:convergence}

\subsection{收敛性定理}
\label{subsec:theorem}

表格型 Q-Learning 的收敛性由 Watkins 和 Dayan（1992）证明\cite{watkins1992q}：

\textbf{定理 3-1（Q-Learning 收敛性）}。对于有限马尔可夫决策过程，若满足以下条件：（1）所有状态 - 动作对被无限次访问；（2）学习率序列满足 Robbins-Monro 条件 $\sum_{t=1}^{\infty} \alpha_t = \infty$，$\sum_{t=1}^{\infty} \alpha_t^2 < \infty$；（3）采用适当的探索策略（如$\varepsilon$-贪婪），则 Q 值序列以概率 1 收敛到最优动作价值函数 $Q^*$。

\subsection{本研究的收敛性保证}
\label{subsec:convergence_guarantee}

在本研究中，收敛性条件分析如下：

\textbf{有限状态空间}。状态空间包含 136 个状态，动作空间包含 4 个动作，Q 表总条目数为 544，满足有限性条件。

\textbf{充分探索}。$\varepsilon$-贪婪策略保证所有状态 - 动作对被充分访问。在 100 轮训练中，每轮 40 步，共 4000 步。544 个状态 - 动作对平均被访问约 7.4 次。考虑到$\varepsilon$-贪婪策略的特性，训练初期的高探索率（$\varepsilon = 0.9$）确保状态空间被广泛探索；随训练进行，Q 值引导的利用过程会使高价值区域被更频繁地访问，保证了关键状态 - 动作对的充分学习。

\textbf{学习率设置}。本研究采用固定学习率 $\alpha = 0.2$，虽不严格满足 Robbins-Monro 条件，但在有限状态空间且充分探索的前提下，可近似收敛到最优解。

因此，本研究算法在理论上保证收敛到最优时间窗选择策略。

\section{与深度 Q 网络的理论对比}
\label{sec:comparison}

\subsection{DQN 方法回顾}
\label{subsec:dqn_review}

深度 Q 网络（Deep Q-Network, DQN）由 Mnih 等人于 2015 年提出\cite{mnih2015human}，使用深度神经网络近似 Q 函数：
\begin{equation}
Q(s, a; \theta) \approx Q^*(s, a)
\end{equation}
其中，$\theta$ 为神经网络参数。DQN 通过最小化以下损失函数训练网络：
\begin{equation}
\mathcal{L}(\theta) = \mathbb{E}_{(s,a,r,s') \sim \mathcal{D}} \left[ \left( r + \gamma \max_{a'} Q(s', a'; \theta^-) - Q(s, a; \theta) \right)^2 \right]
\end{equation}
其中，$\theta^-$ 为目标网络参数，$\mathcal{D}$ 为经验回放缓冲区。

\subsection{理论特性对比}
\label{subsec:theoretical_comparison}

表~\ref{tab:comparison} 从理论层面对比表格型 Q-Learning 与 DQN 的特性。

\begin{table}[H]
\centering
\caption{表格型 Q-Learning 与 DQN 理论特性对比}
\label{tab:comparison}
\begin{tabular}{lcc}
\hline
\textbf{维度} & \textbf{表格型 Q-Learning} & \textbf{DQN} \\
\hline
函数近似 & 无（精确表格表示） & 深度神经网络（通用函数近似） \\
收敛性保证 & 理论保证~\cite{watkins1992q} & 无严格理论保证，经验性收敛~\cite{mnih2015human} \\
状态空间要求 & 有限、可离散化 & 连续或大规模离散 \\
样本复杂度 & $O(|\mathcal{S}||\mathcal{A}|)$ & $O(\text{网络参数量})$ \\
推理复杂度 & $O(1)$ 查表 & $O(\text{网络层数} \times \text{神经元数})$ \\
可解释性 & 高（Q 值透明可视） & 低（黑箱模型） \\
\hline
\end{tabular}
\end{table}

\subsection{适用边界分析}
\label{subsec:applicability}

基于上述理论对比，两种方法的适用边界如下：

\textbf{表格型 Q-Learning 适用场景}：状态空间规模有限（$|\mathcal{S}| < 10^4$）；状态可精确离散化；状态 - 动作对可充分访问；需要严格收敛性保证。

\textbf{DQN 适用场景}：状态空间规模大（$|\mathcal{S}| > 10^5$）或连续；状态表示复杂（如图像输入）；需要跨状态泛化；可接受经验性收敛。

在运动想象时间窗优化任务中，状态空间规模为 136，属于小规模有限状态空间，且状态可精确离散化为时间窗参数。因此，表格型 Q-Learning 是更合适的选择。

\subsection{方法选择依据}
\label{subsec:rationale}

本研究选择表格型 Q-Learning 而非 DQN，基于以下理论考量：

（1）\textbf{状态空间规模}。时间窗参数空间离散化后仅 136 个状态，Q 表总条目数为 544，在现代计算机内存中仅占用约 4 KB 空间（float64 精度）。使用神经网络进行函数近似不仅不必要，还会引入额外的复杂度和不稳定性。

（2）\textbf{收敛性保证}。表格型 Q-Learning 在有限状态空间下具有严格的收敛性保证，而 DQN 的收敛性仅依赖经验验证。对于 BCI 等安全敏感应用，收敛性保证是重要考量因素。

（3）\textbf{样本效率}。表格型方法直接在 Q 表上更新，单步即可传播奖励信号；DQN 需要通过批量梯度下降训练神经网络，样本效率较低。在 BCI 数据稀缺场景下，样本效率是关键指标。

（4）\textbf{可解释性}。表格型 Q-Learning 的 Q 值表可直接可视化和分析，便于理解决策依据和调试优化；DQN 的神经网络权重难以解释，决策过程不透明。

\section{本章小结}

本章提出了一种基于表格型 Q-Learning 的运动想象时间窗优化方法。首先，将时间窗优化问题形式化为马尔可夫决策过程，定义了包含 136 个离散状态的状态空间和包含 4 个边界调整动作的动作空间，采用交叉验证准确率作为奖励信号，并证明了状态转移的确定性特性。其次，详细设计了表格型 Q-Learning 算法，采用$\varepsilon$-贪婪策略平衡探索与利用，通过时序差分学习迭代更新 Q 值表，推理阶段仅需$O(1)$复杂度的查表操作即可确定最优时间窗。再次，从理论上证明了算法的收敛性：在有限状态空间（136 个状态）和$\varepsilon$-贪婪充分探索的前提下，Q 值序列保证收敛到最优动作价值函数$Q^*$。最后，系统对比了表格型 Q-Learning 与 DQN 的理论特性，从状态空间规模、收敛性保证、样本效率和可解释性四个维度论证了本研究选择表格型方法的合理性。

本章方法的创新点在于：（1）\textbf{问题建模简洁性}：将时间窗参数离散化为有限状态空间，避免了深度强化学习的复杂性；（2）\textbf{理论保证严格性}：基于 Watkins 和 Dayan（1992）的收敛性定理，在有限状态空间下具有严格的收敛性保证；（3）\textbf{计算效率优势}：离线训练阶段虽需 4000 次奖励评估（100 episodes × 40 steps），但在线推理阶段仅需单次查表操作，时间复杂度为$O(1)$，完全满足在线 BCI 系统的实时性要求；（4）\textbf{可解释性强}：Q 值表可直接可视化和分析，便于理解决策依据和调试优化。

然而，本章方法仍存在一定局限：（1）状态空间需要人工离散化，可能丢失部分连续信息；（2）奖励计算依赖 5 折交叉验证，训练阶段计算开销较大；（3）仅适用于时间窗优化这一特定任务，泛化能力有限。这些局限性将在第 4 章通过实验进行详细分析，并与多种基线方法（包括固定时间窗、网格搜索、随机搜索和 DQN 等）进行系统对比，全面评估本研究方法的有效性与实用性。
